\documentclass[12pt,a4paper]{article}

\usepackage[utf8]{inputenc}
\usepackage[T1]{fontenc}
\usepackage{lmodern}
\usepackage[a4paper,margin=2.5cm]{geometry}
\usepackage{amsmath,amssymb,amsthm}
\usepackage{hyperref}

\hypersetup{
    colorlinks=true,
    linkcolor=blue,
    urlcolor=blue,
    citecolor=blue,
    pdftitle={Relativistic Phenomenology of ORT},
    pdfauthor={Fedor Kapitanov}
}

\newtheorem{definition}{Definition}
\newtheorem{postulate}{Postulate}

\title{\textbf{Relativistic Phenomenology of Ontological Resolution Theory}\\[0.3em]
\large Curvature as Memory, Expansion as Log Growth, Black Holes as Saturation}

\author{
Fedor Kapitanov\\[0.3em]
\small Independent Researcher, Moscow\\
\small ORCID: 0009-0009-6438-8730\\
\small \texttt{prtyboom@gmail.com}
}

\date{November 2025}

\begin{document}

\maketitle

\begin{abstract}
We present the relativistic phenomenology of Ontological Resolution Theory (ORT), reinterpreting General Relativity in informational terms. The stress-energy tensor $T_{\mu\nu}$ is understood as the density of ontological distinctions, decomposed into active and archival sectors. Cosmological expansion reflects the growth of the universe's holographic log capacity. Black holes are local saturation points of memory, and gravitational waves are coherent perturbations generated by large-scale log updates. This framework preserves the mathematical structure of GR while providing an ontological foundation connecting gravity, information, and thermodynamics.
\end{abstract}

\tableofcontents
\newpage

\section{Introduction}

In a companion paper, we introduced Ontological Resolution Theory (ORT) as a framework unifying quantum mechanics, gravity, dark matter, and cognition through the concept of ontological distinction. The present paper develops the relativistic phenomenology of ORT, showing how standard General Relativity (GR) can be reinterpreted---without modification---as the thermodynamics of distinction.

The goal is not to replace GR but to provide an ontological reading in which:
\begin{itemize}
    \item Curvature encodes the density of recorded distinctions
    \item Expansion reflects the growth of holographic log capacity
    \item Black holes represent local memory saturation
    \item Gravitational waves are coherent log-update perturbations
\end{itemize}

\section{Energy--Momentum as Density of Distinctions}

\subsection{Standard Formulation}

The Einstein field equations relate curvature to matter:
\begin{equation}
    G_{\mu\nu} = \frac{8\pi G}{c^4}\,T_{\mu\nu}
\end{equation}
where $G_{\mu\nu}$ is the Einstein tensor and $T_{\mu\nu}$ is the stress-energy tensor.

\subsection{ORT Interpretation}

Within ORT, $T_{\mu\nu}$ is not a primitive source but an emergent, coarse-grained description of the density of ontological distinctions. We decompose:
\begin{equation}
    T_{\mu\nu}(x) = T_{\mu\nu}^{\rm act}(x) + T_{\mu\nu}^{\rm arch}(x)
\end{equation}
where:
\begin{itemize}
    \item $T_{\mu\nu}^{\rm act}$ encodes the \emph{active sector}: energy-momentum of processes currently producing new distinctions (ordinary matter, radiation, fields)
    \item $T_{\mu\nu}^{\rm arch}$ encodes the \emph{archival sector}: effective energy-momentum of past, fully decohered distinctions---this is dark matter in ORT
\end{itemize}

The Einstein equations become:
\begin{equation}
    G_{\mu\nu}(x) = \frac{8\pi G}{c^4}\bigl(T_{\mu\nu}^{\rm act}(x) + T_{\mu\nu}^{\rm arch}(x)\bigr)
\end{equation}

Formally identical to GR, but ontologically the right-hand side is the macroscopic equation of state for the \emph{log of distinctions}.

\begin{center}
\emph{Curvature is not sourced by ``matter'' in the abstract, but by the integrated history of distinction events.}
\end{center}

\section{Cosmological Expansion as Information Growth}

\subsection{The Holographic Constraint}

By the holographic principle, maximum entropy within a region is proportional to boundary area:
\begin{equation}
    S_{\max} \sim \frac{A}{4\ell_P^2}
\end{equation}

If recorded distinctions grow monotonically:
\begin{equation}
    \frac{dS}{dt} > 0
\end{equation}
then holographic capacity must keep pace:
\begin{equation}
    \frac{dA_{\rm hor}}{dt} \gtrsim 0
\end{equation}

\subsection{Expansion as Log Growth}

In FLRW cosmology, horizon area $A_{\rm hor}(t)$ relates to Hubble radius $R_H \sim c/H$ and scale factor $a(t)$. Persistent distinction production, constrained by holographic bounds, drives horizon area increase and effective expansion.

\begin{itemize}
    \item \textbf{Fact of expansion} ($a(t)$ grows): distinctions accumulate
    \item \textbf{Rate/acceleration} (dark energy): how $\dot{S}(t)$ couples to archival sector and vacuum contributions
\end{itemize}

Accelerated expansion is not mysterious ``anti-gravity'' but the requirement that the holographic screen must enlarge to accommodate growing distinctions.

\begin{center}
\emph{Expansion is the kinematics of keeping the log writable.}
\end{center}

\section{Black Holes as Memory Saturation}

\subsection{Bekenstein-Hawking Entropy}

Black holes carry entropy:
\begin{equation}
    S_{\rm BH} = \frac{k_B c^3 A}{4 G\hbar}
\end{equation}
suggesting maximal entropy objects for given mass and size.

\subsection{ORT Interpretation}

Black holes are \emph{local saturation points of memory capacity}---regions where incoming distinctions are so dense that spacetime cannot expand locally to maintain low-density description.

\begin{itemize}
    \item Horizon area $A_H$ measures local log capacity in Planck bits
    \item Infalling matter records distinctions at/near horizon, increasing $A_H$ and $S_{\rm BH}$
    \item Inside horizon, detailed structure becomes externally inaccessible; description coarse-grains to mass, spin, charge
\end{itemize}

\begin{definition}[Memory Saturation Region]
A black hole is a spacetime region where local information processing switches from ``recording in distributed active degrees of freedom'' to ``archiving in ultra-compressed, effectively inert form.''
\end{definition}

\subsection{Hawking Radiation}

Hawking radiation is information leakage and re-distribution:
\begin{itemize}
    \item Accretion + emission acts as slow, noisy channel converting ultra-compact archive back to delocalized degrees of freedom
    \item Unitarity preserved at fundamental level: distinctions neither created nor destroyed, only re-encoded
\end{itemize}

The information problem rephrases: not ``Does information disappear?'' but ``How is ultra-compressed archival encoding mapped back into outgoing radiation?''

\section{Gravitational Waves as Log-Update Perturbations}

\subsection{Standard View}

GWs are transverse, traceless metric perturbations propagating at $c$, produced by time-varying quadrupole moments.

\subsection{ORT Interpretation}

Gravitational waves are \emph{coherent perturbations in consensus field $\rho_C$ and curvature}, generated by large-scale \emph{rewriting and archiving} operations.

A binary black hole merger:
\begin{itemize}
    \item Fusion of two saturated memory regions into single archival block
    \item Drastic reorganization of local gravitational memory
    \item Global ``re-indexing'' induces oscillatory adjustments detected as GW ringdown
\end{itemize}

Horizon area discreteness at Planck scale:
\begin{equation}
    \Delta A = 8\pi \ell_P^2
\end{equation}
suggests quasi-discrete spectral features---signature of \emph{bitwise updates} to archival sector.

\begin{center}
\emph{Gravitational waves are acoustic waves in the universe's log, generated by large-scale batch updates to gravitational memory.}
\end{center}

\section{Summary: GR as Thermodynamics of Distinction}

\begin{table}[h]
\centering
\begin{tabular}{|l|l|}
\hline
\textbf{GR Concept} & \textbf{ORT Interpretation} \\
\hline
$G_{\mu\nu}$ (curvature) & Memory density \\
$T_{\mu\nu}$ (stress-energy) & Distinction density (active + archival) \\
Expansion $\dot{a} > 0$ & Log capacity growth \\
Dark energy / $\Lambda$ & Write-rate constraint \\
Black holes & Memory saturation regions \\
Hawking radiation & Archive decompression \\
Gravitational waves & Log-update perturbations \\
\hline
\end{tabular}
\caption{ORT reinterpretation of relativistic concepts}
\end{table}

General Relativity is not an isolated geometrical law but the \emph{macroscopic thermodynamic equation of state} of distinction within the Absolute.

\section{Conclusion}

We have shown that GR admits a complete reinterpretation within ORT:
\begin{enumerate}
    \item \textbf{Curvature is memory density}: $G_{\mu\nu}$ encodes how densely distinctions are recorded
    \item \textbf{Expansion is write speed}: $\dot{a}(t)$ reflects holographic capacity allocation
    \item \textbf{Black holes are memory limits}: local saturation with ultra-compressed archival
    \item \textbf{Gravitational waves are log updates}: coherent perturbations from large-scale rewrites
\end{enumerate}

This preserves all predictions of GR while providing ontological foundation connecting gravity, information, and thermodynamics in a unified framework.

\section*{Acknowledgments}

Developed with AI assistance for drafting and consistency checking. Conceptual structure and scientific claims are the author's responsibility.

\begin{thebibliography}{99}

\bibitem{KapitanovORT}
F.~Kapitanov,
\emph{Ontological Resolution Theory: Gravity as Memory, Cognition as Indexing},
Zenodo (2025).

\bibitem{Bousso02}
R.~Bousso,
\emph{The holographic principle},
Rev.\ Mod.\ Phys.\ \textbf{74}, 825 (2002).

\bibitem{Bekenstein73}
J.~D.~Bekenstein,
\emph{Black holes and entropy},
Phys.\ Rev.\ D \textbf{7}, 2333 (1973).

\bibitem{Hawking75}
S.~W.~Hawking,
\emph{Particle creation by black holes},
Commun.\ Math.\ Phys.\ \textbf{43}, 199 (1975).

\bibitem{Jacobson95}
T.~Jacobson,
\emph{Thermodynamics of spacetime: The Einstein equation of state},
Phys.\ Rev.\ Lett.\ \textbf{75}, 1260 (1995).

\end{thebibliography}

\end{document}